\documentclass[11pt,letterpaper]{article}

% Essential Packages
\usepackage[margin=1in]{geometry}
\usepackage{amsmath,amssymb}
\usepackage{graphicx}
\usepackage{circuitikz}
\usepackage{tikz}
\usetikzlibrary{shapes,arrows,positioning,calc}
\usepackage{booktabs}
\usepackage{multirow}
\usepackage{xcolor}
\usepackage{fancyhdr}
\usepackage{tcolorbox}
\usepackage{enumitem}
\usepackage{float}
\usepackage{listings}
\usepackage{hyperref}

% Colors
\definecolor{nlcblue}{RGB}{0,51,102}
\definecolor{alertred}{RGB}{204,0,0}
\definecolor{safgreen}{RGB}{0,128,0}

% Header/Footer
\pagestyle{fancy}
\fancyhead[L]{\textbf{NLC STEM Club UAV Project}}
\fancyhead[R]{\textbf{Safety Protocol}}
\fancyfoot[C]{\thepage}

% Title
\title{\Huge\textbf{Safety Protocol}\\
\Large LiPo Handling and Electrical Safety\\
\large NLC STEM Club - Electrical Systems Team}
\author{February 2026}
\date{}

\begin{document}

\maketitle
\thispagestyle{empty}
\newpage
\section{ELECTRICAL SAFETY PROTOCOL}

\begin{itemize}
  \item Fixed-Wing Environmental Monitoring UAV
  \item NLC STEM Club - Electrical Systems Team
\end{itemize}

\section{MANDATORY READING FOR ALL TEAM MEMBERS}

\section{[WARNING] WARNING [WARNING]}

\begin{itemize}
  \item This project involves high-energy LiPo batteries,
  \item rotating propellers, and high-current electronics.
\end{itemize}

\subsection{Failure to follow safety protocols can result in}

\begin{itemize}
  \item Fire or explosion (LiPo batteries)
  \item Severe lacerations (propeller strikes)
  \item Electrical shock or burns
  \item Equipment damage
  \item Personal injury
\end{itemize}

\section{SAFETY IS NON-NEGOTIABLE.}

\begin{itemize}
  \item When in doubt, STOP and ask for guidance.
\end{itemize}

\begin{itemize}
  \item SECTION 1: LiPo BATTERY SAFETY (HIGHEST PRIORITY)
\end{itemize}

LiPo batteries store significant energy (14.8V $\times$ 2.2Ah = 32.56 Wh). Mishandling can cause\par

thermal runaway, fire, or explosion. LiPo fires burn at >1000$^\circ$C and cannot be extinguished\par

with water.\par

\subsection*{LIPO HANDLING RULES (MANDATORY)}

\subsection*{✓ DO:}

\subsection*{✓ Store batteries in LiPo safety bag at ALL times when not in use}

\subsection*{✓ Charge batteries in LiPo safety bag on non-flammable surface}

\subsection*{✓ Supervise charging at all times (NEVER leave unattended)}

\subsection*{✓ Use only the provided B6 balance charger (4S LiPo mode, 1C rate max)}

\subsection*{✓ Check cell voltages before each use: 3.7V ± 0.1V per cell}

\subsection*{✓ Disconnect battery immediately after use}

\subsection*{✓ Store at 3.8V per cell (storage charge) if not using for >1 week}

\subsection*{✓ Inspect battery before each use: Look for damage, swelling, or deformation}

\subsection*{✓ Transport batteries in LiPo bag, padded container}

\subsection*{✓ Keep fire extinguisher (ABC or CO₂ rated) nearby during charging}

\subsection*{✗ DO NOT:}

\subsection*{✗ NEVER charge unattended or overnight}

\subsection*{✗ NEVER discharge below 3.3V per cell (13.2V total) - causes permanent damage}

\subsection*{✗ NEVER overcharge above 4.2V per cell (16.8V total) - fire risk}

\subsection*{✗ NEVER charge at >1C rate (2.2A max for our 2200mAh battery)}

\subsection*{✗ NEVER puncture, crush, or short-circuit battery terminals}

\subsection*{✗ NEVER store fully charged (fire risk increases over time)}

\subsection*{✗ NEVER use damaged or swollen batteries (dispose properly)}

\subsection*{✗ NEVER expose to water (LiPo + water = dangerous reaction)}

\subsection*{✗ NEVER charge in hot environment (>40°C / 104°F)}

\subsection*{✗ NEVER use damaged charger or wrong battery settings}

\subsection*{BATTERY CHARGING PROCEDURE}

\subsection*{Step 1: Pre-Charge Inspection}

\subsection*{☐ Visually inspect battery (no swelling, no damage to wires or connector)}

\subsection*{☐ Check cell voltages with multimeter: All cells 3.0-4.2V}

\subsection*{☐ If any cell <3.0V or >4.3V → DO NOT CHARGE → Report to lead}

\subsection*{Step 2: Setup}

\subsection*{☐ Place battery in LiPo safety bag}

\subsection*{☐ Place bag on non-flammable surface (metal, ceramic, concrete - NOT wood/carpet)}

\subsection*{☐ Clear area of flammable materials (paper, solvents, aerosols)}

\subsection*{☐ Have fire extinguisher within reach}

\subsection*{Step 3: Charger Configuration}

\subsection*{☐ Connect balance plug to charger (5-pin JST-XH connector)}

\subsection*{☐ Connect main power plug (XT60) to charger}

\subsection*{☐ Set charger to "LiPo Balance Charge" mode}

\subsection*{☐ Set cell count: 4S (14.8V)}

\subsection*{☐ Set charge current: 1.0A (0.5C) or 2.2A max (1C)}

\subsection*{☐ Verify charger displays "4S 14.8V" before starting}

\subsection*{Step 4: Charging}

\subsection*{☐ Press START, charger will check cell count}

\subsection*{☐ Press START again to confirm and begin charging}

\subsection*{☐ REMAIN NEARBY during entire charge cycle (typically 60-120 minutes)}

\subsection*{☐ Monitor battery temperature (should be warm, not hot - <45°C)}

\subsection*{☐ If battery becomes hot (>50°C), smells unusual, or swells → STOP IMMEDIATELY}

\subsection*{Step 5: Post-Charge}

\subsection*{☐ Charger will beep when complete (all cells at 4.20V ± 0.01V)}

\subsection*{☐ Disconnect balance plug first, then main power plug}

\subsection*{☐ Let battery cool for 15 minutes before use or storage}

\subsection*{☐ Label battery with charge date and voltage}

\subsection*{☐ If storing >1 week, discharge to storage voltage (3.8V/cell = 15.2V total)}

\subsection*{LIPO EMERGENCY PROCEDURES}

\subsection*{SWOLLEN BATTERY:}

\subsection*{1. DO NOT CHARGE or USE}

\subsection*{2. Place in LiPo bag, store in safe outdoor location}

\subsection*{3. Discharge to 0V using salt water method (see disposal procedure)}

\subsection*{4. Report to electrical lead for proper disposal}

\subsection*{BATTERY FIRE (Thermal Runaway):}

\subsection*{1. EVACUATE AREA IMMEDIATELY - Alert others}

\subsection*{2. DO NOT USE WATER (makes fire worse)}

\subsection*{3. If safe, use ABC fire extinguisher or sand to smother}

\subsection*{4. Move battery outdoors if possible (use metal tongs, not bare hands)}

\subsection*{5. Call 911 if fire spreads or cannot be controlled}

\subsection*{6. LiPo fires burn for 10-30 minutes - let it burn out in safe area}

\subsection*{BATTERY PUNCTURE/DAMAGE:}

\subsection*{1. Place in LiPo bag immediately}

\subsection*{2. Move to outdoor location away from flammable materials}

\subsection*{3. Monitor for 1 hour (thermal runaway may be delayed)}

\subsection*{4. Discharge and dispose if no fire occurs}

\subsection*{SALT WATER DISCHARGE METHOD (For Disposal):}

\subsection*{1. Prepare salt water solution (5 tablespoons salt per 1 gallon water)}

\subsection*{2. Use non-metallic container (plastic bucket)}

\subsection*{3. Submerge battery completely (weigh down with rock)}

\subsection*{4. Leave submerged for 1 week (check voltage - must be <1V before disposal)}

\subsection*{5. Dispose at battery recycling center (NOT regular trash)}

\section{SECTION 2: PROPELLER SAFETY}

A 10-inch propeller spinning at 8000 RPM has tip speed of \textasciitilde{}400 km/h (\textasciitilde{}250 mph).\par

Contact with spinning propeller causes severe lacerations or amputation.\par

\subsection*{PROPELLER SAFETY RULES (MANDATORY)}

\subsection*{✓ DO:}

\subsection*{✓ Remove propeller for ALL bench testing indoors}

\subsection*{✓ Install propeller ONLY immediately before flight}

\subsection*{✓ Verify propeller securely attached (check spinner nut tight)}

\subsection*{✓ Stand behind propeller plane of rotation when arming motor}

\subsection*{✓ Shout "PROPELLER TEST" before any motor test outdoors}

\subsection*{✓ Establish 10-meter safety perimeter for motor tests}

\subsection*{✓ Inspect propeller before each flight (no cracks, no damage)}

\subsection*{✓ Use folding propeller as designed (blades fold back for water landing)}

\subsection*{✗ DO NOT:}

\subsection*{✗ NEVER test motor with propeller attached indoors}

\subsection*{✗ NEVER reach toward spinning propeller}

\subsection*{✗ NEVER arm motor with people in front of aircraft}

\subsection*{✗ NEVER use damaged or cracked propeller (causes imbalance, shattering risk)}

\subsection*{✗ NEVER overtighten propeller (can crack hub)}

\subsection*{✗ NEVER stand in front of aircraft when battery connected}

\subsection*{MOTOR TEST PROCEDURE}

\subsection*{BENCH TEST (Indoors - Diagnostics only):}

\subsection*{1. Ensure propeller is REMOVED}

\subsection*{2. Secure aircraft to bench (prevent movement)}

\subsection*{3. Connect battery, arm motor}

\subsection*{4. Test throttle response at LOW throttle only (10-20%)}

\subsection*{5. Listen for unusual sounds (grinding, clicking = motor damage)}

\subsection*{6. Check motor temperature after 10-second test (should be cool to touch)}

\subsection*{7. Disarm and disconnect battery immediately after test}

\subsection*{OUTDOOR TEST (With Propeller):}

\subsection*{1. Install propeller (verify secure attachment)}

\subsection*{2. Establish 10-meter safety perimeter with cones/rope}

\subsection*{3. Appoint safety officer to monitor perimeter}

\subsection*{4. Shout "PROPELLER TEST - CLEAR THE AREA"}

\subsection*{5. Wait 10 seconds for confirmation area is clear}

\subsection*{6. Connect battery while standing behind aircraft}

\subsection*{7. Arm motor, gradually increase throttle to 50% max}

\subsection*{8. Test duration: 10 seconds maximum}

\subsection*{9. Throttle to zero, disarm, disconnect battery}

\subsection*{10. Wait for propeller to stop completely before approaching}

\section{SECTION 3: ELECTRICAL SAFETY}

\subsection*{SOLDERING SAFETY}

\subsection*{✓ DO:}

\subsection*{✓ Use soldering iron in well-ventilated area (fume extraction if available)}

\subsection*{✓ Wear safety glasses (solder can splatter)}

\subsection*{✓ Use soldering iron stand (never lay hot iron on bench)}

\subsection*{✓ Set temperature to 350°C (660°F) for lead-free solder}

\subsection*{✓ Wash hands after soldering (lead exposure risk)}

\subsection*{✓ Allow solder joints to cool before handling}

\subsection*{✓ Unplug iron immediately after use}

\subsection*{✓ Keep flammable materials away from soldering area}

\subsection*{✗ DO NOT:}

\subsection*{✗ Touch soldering iron tip (causes severe burns)}

\subsection*{✗ Inhale solder fumes directly (contains flux vapors)}

\subsection*{✗ Solder near LiPo batteries (heat can cause thermal runaway)}

\subsection*{✗ Use wet sponge with excessive water (steam can burn)}

\subsection*{✗ Leave soldering iron unattended while hot}

\subsection*{FIRST AID (Burns):}

\subsection*{1. Run cold water over burn for 10-15 minutes}

\subsection*{2. Do not apply ice directly (tissue damage)}

\subsection*{3. Cover with sterile bandage}

\subsection*{4. Seek medical attention for burns >1 inch or on face/hands}

\subsection*{SHORT CIRCUIT PREVENTION}

\subsection*{A short circuit across 4S LiPo can discharge 286A (burst), causing:}

\subsection*{- Wires to melt or catch fire}

\subsection*{- Battery thermal runaway}

\subsection*{- Burns from hot metal}

\subsection*{PREVENTION:}

\subsection*{✓ Use heat shrink tubing on all solder joints}

\subsection*{✓ Keep tools (screwdrivers, pliers) away from exposed battery terminals}

\subsection*{✓ Disconnect battery during ALL wiring work}

\subsection*{✓ Use XT60 connector covers when battery not in use}

\subsection*{✓ Secure all wires (prevent fraying and short circuits)}

\subsection*{✓ Test continuity BEFORE connecting battery (check for shorts with multimeter)}

\subsection*{EMERGENCY (Short Circuit Occurring):}

\subsection*{1. Disconnect battery IMMEDIATELY (grab insulated connector, not wires)}

\subsection*{2. If fire starts, follow LiPo fire procedures (Section 1)}

\subsection*{3. Inspect all components for damage before re-use}

\subsection*{ELECTROSTATIC DISCHARGE (ESD) PROTECTION}

\subsection*{Flight controller, GPS, and sensors contain sensitive electronics vulnerable to ESD.}

\subsection*{PREVENTION:}

\subsection*{✓ Ground yourself before handling electronics (touch metal bench/ground)}

\subsection*{✓ Use ESD wrist strap if available}

\subsection*{✓ Handle circuit boards by edges (avoid touching components)}

\subsection*{✓ Work on non-conductive surface (wood, plastic mat - NOT metal)}

\subsection*{✓ Store components in anti-static bags when not in use}

\subsection*{✗ Avoid working in low-humidity environments (increases static buildup)}

\subsection*{✗ Avoid wearing synthetic fabrics (wool, polyester generate static)}

\section{SECTION 4: WORKSHOP \& TESTING SAFETY}

\subsection*{GENERAL WORKSHOP SAFETY}

\subsection*{✓ Maintain clean, organized workspace (clutter causes accidents)}

\subsection*{✓ Wear safety glasses during soldering, drilling, cutting}

\subsection*{✓ Keep food and drinks away from electronics (spill risk)}

\subsection*{✓ Label all wires and components clearly (prevents wiring errors)}

\subsection*{✓ Use proper tools for each task (reduces damage and injury risk)}

\subsection*{✓ Report damaged tools or equipment to lead immediately}

\subsection*{✓ Know location of first aid kit and fire extinguisher}

\subsection*{✗ Do not work alone on high-risk tasks (battery charging, motor tests)}

\subsection*{✗ Do not rush (mistakes cause safety hazards)}

\subsection*{✗ Do not bypass safety procedures to save time}

\subsection*{POWER-UP SEQUENCE (EVERY TIME)}

\subsection*{Following this sequence prevents damage and ensures safe operation:}

\subsection*{PRE-POWER CHECKS:}

\subsection*{☐ 1. Propeller REMOVED (if bench testing indoors)}

\subsection*{☐ 2. All connections secure (no loose wires)}

\subsection*{☐ 3. No short circuits (visual inspection of exposed metal)}

\subsection*{☐ 4. Battery voltage checked: 14.0-16.8V total, cells balanced within 0.05V}

\subsection*{☐ 5. Transmitter turned ON (prevents failsafe activation)}

\subsection*{☐ 6. Throttle stick at ZERO position}

\subsection*{POWER-UP:}

\subsection*{☐ 7. Connect battery (listen for beeps from ESC and FC status LED)}

\subsection*{☐ 8. Wait 30 seconds for system initialization}

\subsection*{☐ 9. Check control surface response (move TX sticks)}

\subsection*{☐ 10. Verify GPS lock if needed (wait 2-5 minutes)}

\subsection*{POST-OPERATION:}

\subsection*{☐ 11. Throttle to zero, disarm motor}

\subsection*{☐ 12. Disconnect battery IMMEDIATELY (prevents accidental arming)}

\subsection*{☐ 13. Check battery voltage (should be >14.0V, if lower → recharge)}

\subsection*{☐ 14. Check component temperatures (motor, ESC should be warm, not hot >50°C)}

\subsection*{MULTIMETER SAFETY}

\subsection*{✓ Use correct setting (DC voltage for batteries, continuity for ground checks)}

\subsection*{✓ Start with highest range and work down (prevents overload)}

\subsection*{✓ Red probe to positive (+), black probe to negative (-) or ground}

\subsection*{✓ Do not measure voltage and current simultaneously (damage to meter)}

\subsection*{✗ Do not probe live high-voltage circuits (>50V) without training}

\subsection*{✗ Do not touch metal probe tips while measuring (shock risk)}

\subsection*{✗ Do not measure current in voltage mode (blows fuse)}

\subsection*{COMMON MEASUREMENTS:}

\subsection*{- Battery voltage: DC V setting, 20V range, measure across XT60 terminals}

\subsection*{- Continuity: Continuity mode (beep symbol), checks for 0Ω connections}

\subsection*{- BEC output: DC V setting, 20V range, measure 5V rail to ground}

\subsection*{- Current draw: DC A setting, 10A or 20A range, inline with circuit}

\section{SECTION 5: TEAM ROLES \& RESPONSIBILITIES}

\subsection*{SAFETY OFFICER (Rotating Role Each Meeting)}

\subsection*{Responsibilities:}

\subsection*{✓ Conduct pre-work safety briefing (5 minutes)}

\subsection*{✓ Monitor battery charging (if occurring during meeting)}

\subsection*{✓ Enforce propeller removal during bench testing}

\subsection*{✓ Ensure fire extinguisher and first aid kit accessible}

\subsection*{✓ Stop any unsafe activity immediately (authority to halt work)}

\subsection*{✓ Verify post-work cleanup and battery storage}

\subsection*{Authority: Safety Officer can override any decision if safety concern exists.}

\subsection*{ELECTRICAL LEAD (You)}

\subsection*{Safety Responsibilities:}

\subsection*{✓ Ensure all team members read and sign safety protocol}

\subsection*{✓ Provide safety training before first hands-on work}

\subsection*{✓ Approve all battery charging operations}

\subsection*{✓ Inspect all wiring before first power-up}

\subsection*{✓ Supervise first motor test for each team member}

\subsection*{✓ Maintain safety equipment (fire extinguisher, first aid kit)}

\subsection*{✓ Report all incidents to club advisors immediately}

\subsection*{ALL TEAM MEMBERS}

\subsection*{You are responsible for:}

\subsection*{✓ Reading this entire safety protocol before starting work}

\subsection*{✓ Asking questions if anything is unclear}

\subsection*{✓ Following ALL safety procedures without exception}

\subsection*{✓ Speaking up if you see unsafe behavior (no exceptions for leadership)}

\subsection*{✓ Reporting near-misses or accidents immediately}

\subsection*{✓ Refusing to perform tasks you're not trained for}

\subsection*{"Safety is everyone's responsibility."}

\section{SECTION 6: EMERGENCY CONTACTS \& RESOURCES}

\subsection*{EMERGENCY CONTACTS}

\subsection*{Fire / Medical Emergency:        911}

\subsection*{NLC Campus Security:             [Fill in campus security number]}

\subsection*{Faculty Advisor:                 [Fill in advisor contact]}

\subsection*{Electrical Lead (You):           [Fill in your phone number]}

\subsection*{FIRST AID KIT LOCATION:          [Fill in location in lab/workshop]}

\subsection*{FIRE EXTINGUISHER LOCATION:     [Fill in location]}

\subsection*{NEAREST FIRE ALARM:              [Fill in location]}

\subsection*{EYEWASH STATION:                 [Fill in location if available]}

\subsection*{INCIDENT REPORTING}

\subsection*{ALL incidents must be reported immediately, including:}

\subsection*{- Injuries (burns, cuts, propeller strikes)}

\subsection*{- Equipment damage (component failures, short circuits)}

\subsection*{- Near-misses (situations that COULD have caused injury)}

\subsection*{- Battery incidents (swelling, fire, thermal runaway)}

\subsection*{Reporting Procedure:}

\subsection*{1. Ensure area is safe (turn off power, move to safe location)}

\subsection*{2. Provide first aid if needed (call 911 for serious injuries)}

\subsection*{3. Notify electrical lead and faculty advisor}

\subsection*{4. Document incident (what happened, root cause, corrective action)}

\subsection*{5. Share lessons learned with team (prevent recurrence)}

\subsection*{"Near-misses are learning opportunities. Report them without fear of blame."}

\section{SECTION 7: SAFETY ACKNOWLEDGMENT}

I have read and understood the Electrical Safety Protocol for the NLC STEM Club UAV project.\par

I agree to follow all safety procedures and understand that failure to do so may result in\par

removal from the project for my safety and the safety of others.\par

I understand the following key risks:\par

\begin{itemize}
  \item LiPo batteries can cause fire or explosion if mishandled
  \item Propellers can cause severe injury or amputation
  \item Soldering irons cause severe burns if mishandled
  \item Short circuits can cause fire and equipment damage
  \item I must disconnect the battery during wiring work
\end{itemize}

I commit to:\par

\begin{itemize}
  \item Following the power-up sequence every time
  \item Never charging LiPo batteries unattended
  \item Removing propellers for all indoor bench testing
  \item Using heat shrink tubing on all solder joints
  \item Speaking up if I see unsafe behavior
  \item Asking questions if I'm unsure about a procedure
\end{itemize}

Name (Print): \underline{\hspace{4.95cm}}     Date: \underline{\hspace{2.4cm}}\par

Signature:    \underline{\hspace{4.95cm}}\par

Team Role:    \underline{\hspace{4.95cm}}\par

\begin{itemize}
  \item (Power Team / Flight Electronics / Research Payload / Integration)
\end{itemize}

\subsection{Witnessed by Electrical Lead}

Name:         \underline{\hspace{4.95cm}}     Date: \underline{\hspace{2.4cm}}\par

Signature:    \underline{\hspace{4.95cm}}\par

\section{SAFETY CHECKLISTS FOR COMMON TASKS (QUICK REFERENCE)}

\subsection*{BENCH TESTING CHECKLIST}

\subsection*{Before powering on:}

\subsection*{☐ Propeller removed (if indoors)}

\subsection*{☐ Battery voltage checked (>14.0V)}

\subsection*{☐ All connections secure}

\subsection*{☐ No exposed metal or potential shorts}

\subsection*{☐ Transmitter turned on, throttle at zero}

\subsection*{☐ Safety officer appointed and present}

\subsection*{After testing:}

\subsection*{☐ Battery disconnected immediately}

\subsection*{☐ Battery voltage logged}

\subsection*{☐ Components checked for overheating}

\subsection*{☐ Battery stored in LiPo bag}

\subsection*{SOLDERING CHECKLIST}

\subsection*{☐ Soldering in well-ventilated area}

\subsection*{☐ Safety glasses worn}

\subsection*{☐ Soldering iron in stand when not in use}

\subsection*{☐ No LiPo batteries near soldering area}

\subsection*{☐ Fire extinguisher nearby}

\subsection*{☐ Heat shrink tubing on all solder joints}

\subsection*{☐ Iron unplugged after use}

\subsection*{☐ Hands washed after soldering}

\subsection*{MOTOR/PROPELLER TEST CHECKLIST (OUTDOOR)}

\subsection*{☐ Propeller secure (spinner nut tight)}

\subsection*{☐ 10-meter safety perimeter established}

\subsection*{☐ Safety officer monitoring perimeter}

\subsection*{☐ Verbal warning given ("PROPELLER TEST - CLEAR THE AREA")}

\subsection*{☐ Area confirmed clear (wait 10 seconds after warning)}

\subsection*{☐ Operator standing behind aircraft}

\subsection*{☐ Test duration <10 seconds at <50% throttle}

\subsection*{☐ Battery disconnected after test}

\subsection*{☐ Wait for propeller to stop before approaching}

\section{FINAL SAFETY REMINDER}

This is a REAL aerospace engineering project with REAL risks.\par

Professional aerospace engineers follow strict safety protocols because the consequences\par

of failure are severe. We are holding ourselves to the same standard.\par

Safety is not optional. Safety is not negotiable. Safety is not something we do when\par

convenient.\par

\section{WHEN IN DOUBT, STOP AND ASK.}

There are no stupid questions about safety.\par

It takes courage to speak up about safety concerns - we encourage it.\par

Mistakes are learning opportunities - report them honestly.\par

Your safety and the safety of your teammates is more important than any deadline or\par

demonstration.\par

\section{DOCUMENT INFORMATION}

Version: 1.0\par

Date: February 20, 2026\par

Author: Electrical Lead, NLC STEM Club\par

Status: MANDATORY READING - SIGNATURE REQUIRED BEFORE HANDS-ON WORK\par

Print this document and have ALL team members sign acknowledgment before starting work.\par

\end{document}
